\documentclass{article}
\usepackage[left=2cm, right=2cm, top=2cm, bottom=2cm]{geometry}
\usepackage{graphicx}
\usepackage{amsmath}
\usepackage{amssymb}
\usepackage{amsthm}
\usepackage{fancyhdr}
\usepackage{verbatim}
\usepackage{listings}
\usepackage{xcolor}
\usepackage{pdfpages}
\usepackage{cancel}
\usepackage{physics}
\usepackage{esint}

\lstset{
    language=C++,
    basicstyle=\ttfamily\footnotesize,
    keywordstyle=\color{blue},
    commentstyle=\color{green},
    stringstyle=\color{red},
    numbers=left,
    numberstyle=\tiny\color{gray},
    frame=single,
    breaklines=true,
    captionpos=b,
}


\title{MTH 553: Lecture \# 39}
\author{Cliff Sun}

\newtheorem{theorem}{Theorem}[section]
\newtheorem{lemma}[theorem]{Lemma}
\newtheorem{definition}[theorem]{Definition}
\newtheorem{conjecture}[theorem]{Conjecture}
\newtheorem{proposition}[theorem]{Proposition}
\newtheorem{corollary}[theorem]{Corollary}
\newtheorem{one minute paper}[theorem]{One Minute Paper}

\pagestyle{fancy}
\lhead{\textbf{\thepage}\ \ \nouppercase{\rightmark}}
\chead{MTH 553: Lecture \# 39}
\rhead{Cliff Sun}

\begin{document}

\maketitle

\section*{Lecture Span}
\begin{enumerate}
    \item Continuation of trajectory
\end{enumerate}

Let 

\begin{equation}
    E(v,w) = \int_{0}^{v}f(z)dz + \frac{1}{2}w^2
\end{equation}

We showed that 

\begin{equation}
    \frac{d}{ds}E(v(s), w(s)) = cw^2(s) \geq 0
\end{equation}

This means that the energy increases along trajectories. We note that 

\begin{equation}
    E(0,0) < E(1,0)
\end{equation}

Since 

\begin{equation}
    \int_{0}^{1}f(v)dv > 0
\end{equation}

For large $c$, the trajectory cuts through the level sets (lines where the energy is constant) quickly. That is, the energy increases rapidly. Here, then the trajectories do not meet. 

For small $c$, the energy will change slowly, but then the trajectory will not meet (undershot). 

Therefore, there must be some in-between value of $c$ of which these trajectories will meet. 

\section*{Schr\"odinger's equation on $\mathbb{R}^n$}

Only difference between SE and heat equation is "i". Recall, we found the fundamental solution to the heat equation by Fourier Transform. 

\begin{equation}
    u_t = \triangle u
\end{equation}

With 

\begin{equation}
    u(x,t) = \frac{1}{(4\pi t)^{n/2}}\int_{\mathbb{R}^n}\exp(-|x-y|^2/4t)g(y)dy
\end{equation}

Then Schr\"odinger's equation 

\begin{equation}
    iu_t + \triangle u = 0
\end{equation}

On $\mathbb{R}^n$. With $u(x,t) \in \mathbb{C}$, or is a complex function. Just the heat equation but with $t$ replaced by $it$. Then 

\begin{equation}
    \frac{\partial u}{\partial (it)} = \triangle u \iff \frac{1}{i}\partial_t u - \triangle u =0
\end{equation}

Multiply by $i^2 = -1$ to obtain SE. So then replacing $t$ with $it$ should obtain the fundamental solution of Schr\"odinger's equation. 

\begin{equation}
    u(x,t) = \frac{1}{\left( 4\pi it \right)^{n/2}}\int_{\mathbb{R}^n}\exp(i|x-y|^2/4t)g(y)dy
\end{equation}

Should solve Schr\"odinger's equation with $u = g(x)$ at $t=0$. For $x \in \mathbb{R}$ and for all $t \in \mathbb{R}$ and $t \neq 0$. You can check (9) directly by differentiating through the integral. We also assume that 
$g \in L^1$ and $|y^2|g(y) \in L^1$. That is, $g$ is integrable. Does $u(x,t) \rightarrow g(x)$ at $t \rightarrow 0$? Yes, using the method of stationary phases. 

SE is time reversible. That is $u(x,-t)$ also solves (SE). For the Schr\"odinger on the bounded domain $\Omega \subset \mathbb{R}^n$ with Dirichlet boundary conditions, then the solution 
is written in terms of eigenfunctions $\{\phi_k\}$. 

\begin{equation}
    u(x,t) = \sum_{k}c_k \exp(-i \lambda_k t)\phi_k(x)
\end{equation}



\end{document}